\begin{figure*}[t]
\centering
\begin{tikzpicture}[caja/.style={draw=black!55, rounded corners=2pt, align=center, text width=1.92cm,
  minimum height=1.75cm, font=\scriptsize, fill=black!3, inner sep=3pt},
  clave/.style={caja, draw=black, line width=0.8pt},
  flecha/.style={-{Stealth[length=4pt]}, black!70, line width=0.6pt}, node distance=0.36cm]
\node[caja] (f) {\textbf{Fuentes públicas}\\[2pt] SIAP 2003 a 2025, CHIRPS, CHIRTS-ERA5, Copernicus DEM e INEGI};
\node[caja, right=of f] (v) {\textbf{Variables por municipio}\\[2pt] promedio sobre la agricultura de temporal, ventanas del ciclo};
\node[caja, right=of v] (t) {\textbf{Tabla municipio-año}\\[2pt] 1\,942 filas de 100 municipios, nivel llevado a 2025};
\node[clave, right=of t] (val) {\textbf{Validación}\\[2pt] espacial con exclusión de 0 a 80 km, por DDR, sin repetidos y por año};
\node[caja, right=of val] (m) {\textbf{Modelos}\\[2pt] bases sin ambiente, simples, kriging de residuos y 14 variables};
\node[clave, right=of m] (r) {\textbf{Regla de decisión}\\[2pt] gana en 7 de 10 y $p \le 0.05$ contra el kriging};
\node[caja, right=of r] (c) {\textbf{Capa de riesgo}\\[2pt] 268 municipios con p10 a p90, probabilidad bajo 2\,t/ha y confianza};
\foreach \a/\b in {f/v, v/t, t/val, val/m, m/r, r/c} \draw[flecha] (\a) -- (\b);
\end{tikzpicture}
\caption{Proceso del modelo. Los recuadros con borde grueso se fijaron antes de correr la versión final de los
experimentos.}
\label{fig:proceso}
\end{figure*}
